% ------------------------------------------------------------------
% Bloques de código.
% Fondo gris claro, numeración lateral y caption superior con doble línea.
% ------------------------------------------------------------------

\renewcommand{\lstlistingname}{Código}
\renewcommand{\lstlistlistingname}{Índice de códigos}

\DeclareCaptionFormat{CodeCaptionRule}{%
  \hrule height 0.5pt\relax
  \vspace{3pt}
  #1#2#3\par
  \vspace{3pt}
  \hrule height 0.5pt\relax
}

\captionsetup[lstlisting]{
  format=CodeCaptionRule,
  position=top,
  font=footnotesize,
  labelfont=bf,
  singlelinecheck=false
}

\lstdefinestyle{BaseCode}{
  backgroundcolor=\color{CodeBg},
  basicstyle=\small\ttfamily,
  commentstyle=\color{CodeComment}\itshape,
  keywordstyle=\color{CodeKeyword}\bfseries,
  stringstyle=\color{CodeString},
  numbers=left,
  numberstyle=\scriptsize\color{CodeComment},
  numbersep=8pt,
  stepnumber=1,
  breaklines=true,
  breakatwhitespace=false,
  showstringspaces=false,
  keepspaces=true,
  columns=fullflexible,
  frame=single,
  framerule=0.4pt,
  rulecolor=\color{CodeBorder},
  framesep=6pt,
  xleftmargin=2.2em,
  framexleftmargin=2.2em,
  xrightmargin=0pt,
  linewidth=\dimexpr\linewidth-2.2em\relax,
  aboveskip=10pt,
  belowskip=10pt,
  tabsize=2,
  captionpos=t
}

\lstdefinelanguage{JavaScript}{
  morekeywords={await,async,break,const,continue,default,export,false,from,function,if,import,let,new,null,return,switch,this,true,try,var,while},
  sensitive=true,
  morecomment=[l]{//},
  morecomment=[s]{/*}{*/},
  morestring=[b]",
  morestring=[b]'
}

\lstdefinelanguage{JSON}{
  morestring=[b]",
  morestring=[b]',
  morecomment=[l]{//},
  morecomment=[s]{/*}{*/},
  morekeywords={true,false,null},
  sensitive=true,
  alsoletter={},
  alsodigit={},
  literate=
    {á}{{\'a}}1 {é}{{\'e}}1 {í}{{\'i}}1 {ó}{{\'o}}1 {ú}{{\'u}}1
    {Á}{{\'A}}1 {É}{{\'E}}1 {Í}{{\'I}}1 {Ó}{{\'O}}1 {Ú}{{\'U}}1
}

% Resaltar específicamente propiedades de JSON
\lstdefinelanguage{JSONExtended}{
  morestring=[b]",
  morestring=[b]',
  morecomment=[l]{//},
  morecomment=[s]{/*}{*/},
  morekeywords={true,false,null},
  sensitive=true,
  alsoletter={},
  alsodigit={},
  % Resaltar nombres de propiedades (experimental)
  literate=
    *{true}{{{\color{CodeKeyword}true}}}4
    {false}{{{\color{CodeKeyword}false}}}5
    {null}{{{\color{CodeKeyword}null}}}4
}

% Entorno para bloques JSON: mismo contador que \lstlistingname
% pero con el prefijo "JSON" en el caption en vez de "Código".
\lstnewenvironment{jsonlisting}[1][]{%
  \lstset{style=json, numbers=none}%
  \lstset{#1}%
  \renewcommand{\lstlistingname}{Código JSON}%
}{%
  \renewcommand{\lstlistingname}{Código}%
}

\lstdefinestyle{ComandoInline}{
  style=BaseCode,
  frame=none,
  numbers=none,
  backgroundcolor=\color{CodeBg},
  basicstyle=\small\ttfamily,
  aboveskip=6pt,
  belowskip=6pt,
  columns=fullflexible,
  breaklines=true,
  breakatwhitespace=false,
}

\NewDocumentCommand{\comando}{v}{%
  \begin{center}
  \lstinline[style=ComandoInline]{#1}
  \end{center}
}


\lstdefinestyle{python}{style=BaseCode,language=Python}
\lstdefinestyle{javascript}{style=BaseCode,language=JavaScript}
\lstdefinestyle{bash}{style=BaseCode,language=bash}
\lstdefinestyle{json}{style=BaseCode,language=JSON}
\lstdefinestyle{text}{style=BaseCode,language={}}
\lstdefinestyle{yaml}{style=BaseCode,language={}}
\lstdefinestyle{sql}{style=BaseCode,language=SQL}
\lstdefinestyle{json}{style=BaseCode,language=JSON}
\lstdefinestyle{jsonextended}{style=BaseCode,language=JSONExtended}

\lstset{style=BaseCode}
